\subsection{The Prime Radical; Prime and Semiprime Rings}
\setcounter{subsection}{1}
skipping chapter/section 3 cz we didnt cover it

\begin{exr}[title=Proposition 10.2] %NOTE
	Prove that the following are equivalent for
	a proper ideal $P$ in a ring $R$. \vspace{-0.2in}
	\begin{enumerate}[i)]
		\item $P$ is a prime ideal, \textit{i.e.} if $I,J$ are two ideals of $R$ and
			$IJ\subseteq P$, then $I\subseteq P$ or $J\subseteq P$.
		\item If $(a)(b)\subseteq P$, then $a\in P$ or $b\in P$.
		\item If $aRb\subseteq P$, then $a\in P$ or $b\in P$.
		\item If $I,J$ are two left ideals of $R$ and $IJ\subseteq P$,
			then $I\subseteq P$ or $J\subseteq P$.
		\item If $I,J$ are two right ideals of $R$ and $IJ\subseteq P$,
			then $I\subseteq P$ or $J\subseteq P$.
	\end{enumerate}
\end{exr}


\nq
\begin{exr}[title=Proposition 10.9] %NOTE
	Prove that the following are equivalent for
	a proper ideal $P$ in a ring $R$. \vspace{-0.2in}
	\begin{enumerate}[i)]
		\item $P$ is a semiprime ideal, \textit{i.e.} if $I$ is an ideal of $R$ and
			$I^2\subseteq P$, then $I\subseteq P$.
		\item If $(a)^2\subseteq P$, then $(a)\subseteq P$.
		\item If $aRa\subseteq P$, then $a\in P$.
		\item If $I$ is a left ideal of $R$ and $I^2\subseteq P$,
			then $I\subseteq P$.
		\item If $I$ is a right ideal of $R$ and $I^2\subseteq P$,
			then $I\subseteq P$.
	\end{enumerate}
\end{exr}


\nq
\begin{exr}[title=Theorem 10.11] %NOTE
	Prove that the following are equivalent for a ring $R$. \vspace{-0.2in}
	\begin{enumerate}[i)]
		\item $R$ is a semiprime ring, meaning that $\{0\}$ is a semiprime ideal of $R$.
		\item $\{0\}$ is an intersection of prime ideals of $R$.
		\item $R$ has no nonzero nilpotent ideals.
		\item $R$ has no nonzero nilpotent \textit{left} ideals.
	\end{enumerate}
	Deduce that if $R$ is a ring in which there are nonzero nilpotents, then $R$ is a \\
	semiprime ring.
\end{exr}


\nq
\begin{exr}[title=Theorem 10.24] %NOTE
	Let $R$ be a ring. Prove that $R$ is a semisimple ring if and only if
	$R$ is semiprime and left artinian.
\end{exr}


\nq
\begin{exr}[title=Exercise 10.0]
	Let $R$ be a prime ring (\textit{i.e.} $\{0\}$ is a prime ideal).
	Prove that $Z(R)$ is an integral domain.
\end{exr} \

\begin{pf}
	Let $a,b\in Z(R)$ such that $ab=0$. Consider any element $x\in aRb$.
	Then for some $r\in R$, we have that:
	\begin{equation*}
		x \ = \ arb \ = \ abr \ = \ 0
	\end{equation*}
	Thus $xRy\subseteq(0)$.
	Since $R$ is prime (i.e. (0) is a prime ideal), we must have that $a=0$ or
	$b=0$ as needed.
\end{pf}


\nq
\begin{exr}[title=Exercise 10.1]
	Show that if $R$ is a semiprime ring, then
	$Z(R)$ is a reduced ring (\textit{i.e.} no nonzero nilpotents).
\end{exr} \

\begin{pf}
	For a fixed $a\in Z(R)$, first consider the case that $a^{2}=0$.
	Since $a\in Z(R)$, we have that $aRa = a^{2}R = 0$, so $aRa\subseteq(0)$.
	Since $R$ (and thus (0)) is semiprime, then we must have that $a=0$.

	Now, suppose $a^{n}=0$ for some arbitrary $n \in\bN$ with $n>1$. Consider:
	\begin{gather*}
		(a^{n/2})^{2} \qquad \trm{if } n \trm{ is even} \\
		\left(a^{\frac{n+1}{2}}\right)^{2} \qquad \trm{if } n \trm{ is odd}
	\end{gather*}
	Applying the above, we see that $a^{m}=0$ for some integer $0<m<n$.
	By infinite descent, we therefore must have that $a=0$ as needed.
\end{pf}


\nq
\begin{exr}[title=Exercise 10.3]
	Let $R$ be a ring.
	Prove that $R$ is a domain if and only if $R$ is prime and reduced.
\end{exr} \

\begin{pf}
	First, suppose $R$ is a domain. If $a^{n}=0$, then:
	\begin{equation*}
		a^{n} \ = \ aa^{n-1} \ = \ 0
	\end{equation*}
	Thus, we necessarily have $a=0$, so $R$ is reduced.
	Now, for any $a,b\in R$, suppose $aRb\subseteq(0)$.
	Then for any $r\neq0$, notice that:
	\begin{equation*}
		arb = a(rb) = 0 \ \implies \ a = 0 \trm{ or } rb = 0
		\ \implies \ a = 0 \trm{ or } b = 0
	\end{equation*}
	Next, suppose $R$ is a prime reduced ring and suppose $ab=0$ for some $a,b\in R$.
	Notice that:
	\begin{equation*}
		(ba)^{2} \ = \ (ba)(ba) \ = \ b(ab)a \ = \ 0
	\end{equation*}
	Since $R$ is reduced, this means $ba=0$ as well.
	Therefore, for any $r\in R$, we have:
	\begin{equation*}
		(arb)^{2} \ = \ (arb)(arb) \ = \ ar(ba)rb \ = \ 0
	\end{equation*}
	Again, since $R$ is reduced, this means $arb=0$, so $aRb\subseteq(0)$.
	Since $R$ is prime, then we must have $a=0$ or $b=0$ as needed.
\end{pf}


\nq
\begin{exr}[title=Exercise 10.4] %TODO
	Prove that if $R$ is left artinian, then every prime ideal of $R$ is maximal.
\end{exr}


\nq
\begin{exr}[title=Exercise 10.8] %TODO
	Let $R$ be a ring. Prove that $R$ is a semiprime ring if and only if
	$IJ=0$ implies $I\cap J=0$ for all ideals $I,J$ in $R$.
\end{exr}
